\chapter*{Введение}                         % Заголовок
\addcontentsline{toc}{chapter}{Введение}    % Добавляем его в оглавление

\newcommand{\actuality}{}
\newcommand{\progress}{}
\newcommand{\aim}{{\text\aimTXT}}
\newcommand{\tasks}{\text{\tasksTXT}}
\newcommand{\novelty}{\text{\noveltyTXT}}
\newcommand{\influence}{\text{\influenceTXT}}
\newcommand{\methods}{\text{\methodsTXT}}
\newcommand{\defpositions}{\text{\defpositionsTXT}}
\newcommand{\reliability}{\text{\reliabilityTXT}}
\newcommand{\probation}{\text{\probationTXT}}
\newcommand{\contribution}{\text{\contributionTXT}}
\newcommand{\publications}{\text{\publicationsTXT}}


{\actuality} 

Компьютерное моделирование широко используется для решения сложных задач в инженерных областях и науке.
%
Этот инструмент особенно важен для моделирование движения жидкости,
так как зачастую аналитического решения возникающих в данной области задач не существует,
и в подобных случаях численные методы являются незаменимыми.
%
В области моделирования гидродинамики ведущую роль играют сеточные методы, такие как 
метод конечных разностей (Finite Difference Method, FDM), 
метод конечных объёмов (Finite Volume Method, FVM) и
метод конечных элементов (Finite Element Method, FEM).
%
Однако у них есть недостатки, связанные со сложностью обработки свободной поверхности жидкости, сложной геометрии, движущихся поверхностей, а также многофазных течений.


В качестве альтернативного подхода всё чаще используются бессеточные методы, 
хорошо подходящие для решения подобных задач, такие как метод сглаженных частиц (Smoothed Particle Hydrodynamics, SPH).
%
Метод SPH был изначально разработан для решения задач астрофизики~\cite{Lucy1977,Monaghan1977,Monaghan1985_astro}, 
и далее стал применяться в других областях, таких как 
газодинамика~\cite{Monaghan1983,Monaghan1985,Monaghan1988,Monaghan1989,Monaghan1992,Monaghan2005}
--- например, для расчёта аэродинамики зданий со сложной геометрией~\cite{Khopeskov2019};
гидродинамика~\cite{Monaghan1994,Monaghan1999,Monaghan2015}
--- для решения множества различных задач, включающих моделирование волн на поверхности воды: 
их распространение~\cite{Dalrymple2001,Dao2013,Dominguez2019,Trimulyono2019} 
и взаимодействие с конструкциями~\cite{Li2004,Ni2013,Lu2021},
многофазные течения~\cite{Monaghan1995,Dao2011,Monaghan2015_Kwon,Wang2016,Wu2017},
и разрушение плотин~\cite{Afanasiev2009,Dal2017,Luo2017};
механика деформируемого твёрдого тела~\cite{Libersky1991,Libersky1993,Gray2001,Gerasimov2010}.
%
У данного метода есть свои недостатки. 
Связаны они со сложностями в обработке некоторых граничных условий, 
в частности со снижением точности из-за нерегулярного распределения частиц~\cite{Chen2000,Liu2006,Korzilius2017}, 
а также неустойчивостью при растяжении~\cite{Swegle1995,Monaghan2000}.

% {\progress}
% Этот раздел должен быть отдельным структурным элементом по
% ГОСТ, но он, как правило, включается в описание актуальности
% темы. Нужен он отдельным структурынм элемементом или нет ---
% смотрите другие диссертации вашего совета, скорее всего не нужен.

{\aim} данной работы является разработка кроссплатформенного программного комплекса для моделирования и анализа сложных гидродинамических процессов, 
с применением суперкомпьютерных вычислительных технологий в инженерных областях.

Для~достижения поставленной цели необходимо было решить следующие {\tasks}:
\begin{enumerate}[beginpenalty=10000] % https://tex.stackexchange.com/a/476052/104425
  \item Исследовать существующие реализации метода SPH и их отличительные особенности.
  \item Подготовить модульную архитектуру проекта для кроссплатформенной реализации в условиях применения различных технологических решения.
  \item Разработать вычислительные модули:
  \begin{enumerate}
    \item \verb|RRSPH_OMP| --- многопоточная реализация на CPU (с применением технологии OpenMP).
    \item \verb|RRSPH_CL| --- гетерогенные вычисления (с применением технологии OpenCL).
  \end{enumerate}
  \item Разработать средства визуализации и обработки:
  \begin{enumerate}
    \item \verb|SPH2D_Drawer| --- средство визуализации двумерных расчётов.
    \item \verb|WaterProfile| --- средство определения профиля уровня жидкости во времени и пространстве.
    \item \verb|FuncAtPoint| --- средство определения значения произвольной функции от времени в указанной точке.
  \end{enumerate}
  \item Разработать скрипты для подготовки входных данных экспериментов.
  \item Обеспечить сборку с использованием популярных инструментов на ОС Windows и Linux, а также работу на данных системах.
  \item Провести тестирование и верификацию программного комплекса:
  \begin{enumerate}
    \item Провести тестирование на эталонных задачах гидродинамики.
    \item Оценить эффективность распараллеливания и производительность различных реализаций вычислительных модулей.
    \item Оценить точность моделирования относительно аналитических и опубликованных численных решений.
  \end{enumerate}
\end{enumerate}


{\novelty}
\begin{enumerate}[beginpenalty=10000] % https://tex.stackexchange.com/a/476052/104425
  \item Впервые \ldots
  \item Впервые \ldots
  \item Было выполнено оригинальное исследование \ldots
\end{enumerate}

{\influence} \ldots

% {\methods} \ldots

% {\defpositions}
% \begin{enumerate}[beginpenalty=10000] % https://tex.stackexchange.com/a/476052/104425
%   \item Первое положение
%   \item Второе положение
%   \item Третье положение
%   \item Четвертое положение
% \end{enumerate}
% В папке Documents можно ознакомиться с решением совета из Томского~ГУ
% (в~файле \verb+Def_positions.pdf+), где обоснованно даются рекомендации
% по~формулировкам защищаемых положений.

% {\reliability} полученных результатов обеспечивается \ldots \ Результаты находятся в соответствии с результатами, полученными другими авторами.


%{\probation}
%{\contribution} Автор принимал активное участие \ldots
По результатам выполнения данной работы имеются апробация и публикации.
\begin{enumerate}[beginpenalty=10000] % https://tex.stackexchange.com/a/476052/104425
    \item Доклад на XXIX региональной конференции молодых учёных и исследователей Волгоградской области.
    \item Доклад на XXVIII региональной конференции молодых учёных и исследователей Волгоградской области.
    \item Доклад на XXVII региональной конференции молодых учёных и исследователей Волгоградской области.
    \item Доклад в Научной сессии ВолГУ 2023.
    \item Статья <<Simulation of Free-Surface Fluid Dynamics: Parallelization for GPUs>> в журнале Lecture Notes in Computer Science 
    (в соавторстве с Савиным Е.С., Храповым С.С., Хоперсковым А.В.)~\cite{Savin2023}
\end{enumerate}


% \ifnumequal{\value{bibliosel}}{0}
% {%%% Встроенная реализация с загрузкой файла через движок bibtex8. (При желании, внутри можно использовать обычные ссылки, наподобие `\cite{vakbib1,vakbib2}`).
%     {\publications} Основные результаты по теме диссертации изложены
%     в~XX~печатных изданиях,
%     X из которых изданы в журналах, рекомендованных ВАК,
%     X "--- в тезисах докладов.
% }%
% {%%% Реализация пакетом biblatex через движок biber
%     \begin{refsection}[bl-author, bl-registered]
%         % Это refsection=1.
%         % Процитированные здесь работы:
%         %  * подсчитываются, для автоматического составления фразы "Основные результаты ..."
%         %  * попадают в авторскую библиографию, при usefootcite==0 и стиле `\insertbiblioauthor` или `\insertbiblioauthorgrouped`
%         %  * нумеруются там в зависимости от порядка команд `\printbibliography` в этом разделе.
%         %  * при использовании `\insertbiblioauthorgrouped`, порядок команд `\printbibliography` в нём должен быть тем же (см. biblio/biblatex.tex)
%         %
%         % Невидимый библиографический список для подсчёта количества публикаций:
%         \phantom{\printbibliography[heading=nobibheading, section=1, env=countauthorvak,          keyword=biblioauthorvak]%
%         \printbibliography[heading=nobibheading, section=1, env=countauthorwos,          keyword=biblioauthorwos]%
%         \printbibliography[heading=nobibheading, section=1, env=countauthorscopus,       keyword=biblioauthorscopus]%
%         \printbibliography[heading=nobibheading, section=1, env=countauthorconf,         keyword=biblioauthorconf]%
%         \printbibliography[heading=nobibheading, section=1, env=countauthorother,        keyword=biblioauthorother]%
%         \printbibliography[heading=nobibheading, section=1, env=countregistered,         keyword=biblioregistered]%
%         \printbibliography[heading=nobibheading, section=1, env=countauthorpatent,       keyword=biblioauthorpatent]%
%         \printbibliography[heading=nobibheading, section=1, env=countauthorprogram,      keyword=biblioauthorprogram]%
%         \printbibliography[heading=nobibheading, section=1, env=countauthor,             keyword=biblioauthor]%
%         \printbibliography[heading=nobibheading, section=1, env=countauthorvakscopuswos, filter=vakscopuswos]%
%         \printbibliography[heading=nobibheading, section=1, env=countauthorscopuswos,    filter=scopuswos]}%
%         %
%         \nocite{*}%
%         %
%         {\publications} Основные результаты по теме диссертации изложены в~\arabic{citeauthor}~печатных изданиях,
%         \arabic{citeauthorvak} из которых изданы в журналах, рекомендованных ВАК%
%         \ifnum \value{citeauthorscopuswos}>0%
%             , \arabic{citeauthorscopuswos} "--- в~периодических научных журналах, индексируемых Web of~Science и Scopus%
%         \fi%
%         \ifnum \value{citeauthorconf}>0%
%             , \arabic{citeauthorconf} "--- в~тезисах докладов.
%         \else%
%             .
%         \fi%
%         \ifnum \value{citeregistered}=1%
%             \ifnum \value{citeauthorpatent}=1%
%                 Зарегистрирован \arabic{citeauthorpatent} патент.
%             \fi%
%             \ifnum \value{citeauthorprogram}=1%
%                 Зарегистрирована \arabic{citeauthorprogram} программа для ЭВМ.
%             \fi%
%         \fi%
%         \ifnum \value{citeregistered}>1%
%             Зарегистрированы\ %
%             \ifnum \value{citeauthorpatent}>0%
%             \formbytotal{citeauthorpatent}{патент}{}{а}{}%
%             \ifnum \value{citeauthorprogram}=0 . \else \ и~\fi%
%             \fi%
%             \ifnum \value{citeauthorprogram}>0%
%             \formbytotal{citeauthorprogram}{программ}{а}{ы}{} для ЭВМ.
%             \fi%
%         \fi%
%         % К публикациям, в которых излагаются основные научные результаты диссертации на соискание учёной
%         % степени, в рецензируемых изданиях приравниваются патенты на изобретения, патенты (свидетельства) на
%         % полезную модель, патенты на промышленный образец, патенты на селекционные достижения, свидетельства
%         % на программу для электронных вычислительных машин, базу данных, топологию интегральных микросхем,
%         % зарегистрированные в установленном порядке.(в ред. Постановления Правительства РФ от 21.04.2016 N 335)
%     \end{refsection}%
%     \begin{refsection}[bl-author, bl-registered]
%         % Это refsection=2.
%         % Процитированные здесь работы:
%         %  * попадают в авторскую библиографию, при usefootcite==0 и стиле `\insertbiblioauthorimportant`.
%         %  * ни на что не влияют в противном случае
%         \nocite{vakbib2}%vak
%         \nocite{patbib1}%patent
%         \nocite{progbib1}%program
%         \nocite{bib1}%other
%         \nocite{confbib1}%conf
%     \end{refsection}%
%         %
%         % Всё, что вне этих двух refsection, это refsection=0,
%         %  * для диссертации - это нормальные ссылки, попадающие в обычную библиографию
%         %  * для автореферата:
%         %     * при usefootcite==0, ссылка корректно сработает только для источника из `external.bib`. Для своих работ --- напечатает "[0]" (и даже Warning не вылезет).
%         %     * при usefootcite==1, ссылка сработает нормально. В авторской библиографии будут только процитированные в refsection=0 работы.
% }

% При использовании пакета \verb!biblatex! будут подсчитаны все работы, добавленные
% в файл \verb!biblio/author.bib!. Для правильного подсчёта работ в~различных
% системах цитирования требуется использовать поля:
% \begin{itemize}
%         \item \texttt{authorvak} если публикация индексирована ВАК,
%         \item \texttt{authorscopus} если публикация индексирована Scopus,
%         \item \texttt{authorwos} если публикация индексирована Web of Science,
%         \item \texttt{authorconf} для докладов конференций,
%         \item \texttt{authorpatent} для патентов,
%         \item \texttt{authorprogram} для зарегистрированных программ для ЭВМ,
%         \item \texttt{authorother} для других публикаций.
% \end{itemize}
% Для подсчёта используются счётчики:
% \begin{itemize}
%         \item \texttt{citeauthorvak} для работ, индексируемых ВАК,
%         \item \texttt{citeauthorscopus} для работ, индексируемых Scopus,
%         \item \texttt{citeauthorwos} для работ, индексируемых Web of Science,
%         \item \texttt{citeauthorvakscopuswos} для работ, индексируемых одной из трёх баз,
%         \item \texttt{citeauthorscopuswos} для работ, индексируемых Scopus или Web of~Science,
%         \item \texttt{citeauthorconf} для докладов на конференциях,
%         \item \texttt{citeauthorother} для остальных работ,
%         \item \texttt{citeauthorpatent} для патентов,
%         \item \texttt{citeauthorprogram} для зарегистрированных программ для ЭВМ,
%         \item \texttt{citeauthor} для суммарного количества работ.
% \end{itemize}
% % Счётчик \texttt{citeexternal} используется для подсчёта процитированных публикаций;
% % \texttt{citeregistered} "--- для подсчёта суммарного количества патентов и программ для ЭВМ.

% Для добавления в список публикаций автора работ, которые не были процитированы в
% автореферате, требуется их~перечислить с использованием команды \verb!\nocite! в
% \verb!Synopsis/content.tex!.
 % Характеристика работы по структуре во введении и в автореферате не отличается (ГОСТ Р 7.0.11, пункты 5.3.1 и 9.2.1), потому её загружаем из одного и того же внешнего файла, предварительно задав форму выделения некоторым параметрам

% \text{Объем и структура работы.} Диссертация состоит из~введения,
% \formbytotal{totalchapter}{глав}{ы}{}{},
% заключения и
% \formbytotal{totalappendix}{приложен}{ия}{ий}{}.
%% на случай ошибок оставляю исходный кусок на месте, закомментированным
%Полный объём диссертации составляет  \ref*{TotPages}~страницу
%с~\totalfigures{}~рисунками и~\totaltables{}~таблицами. Список литературы
%содержит \total{citenum}~наименований.
%

В первой главе рассмотрены уравнения Навье-Стокса и их дискретизация в рамках метода SPH.
%
Особое внимание уделено обеспечения несжимаемости жидкости и алгоритмам временного интегрирования,
включая критерии их устойчивости и выбор шага по времени.
%
Описаны методы задания граничных условий, критически важные для задач со сложной геометрией.
%
Описаны подходы к генерации и поглощению волн на поверхности жидкости, 
необходимые для моделирования гидродинамических процессов.

% Представлены различные методы интегрирования по времени, их устойчивость и связанные с ней формулы для выбора шага по времени.
% От выбора граничных условий часто многое зависит, особенно при сложной геометрии, 
% либо когда много частиц находится вблизи границ, поэтому дальше рассматриваются методы обработки граничных условий в SPH.
% При моделировании воды часто требуется генерировать волны, и разрабатываемое ПО должно быть в состоянии это делать. 
% Поэтому рассматриваются различные методы генерации и погощения волн на воде.

Во второй главе представлена архитектура программного комплекса для суперкомпьютерного моделирования методом SPH.
%
Подробно описаны вычислительные модули, средства визуализации, постобработки и вспомогательные скрипты:
их возможности, особенности и реализация.
%
Описаны потоки данных между компонентами системы, а также ключевые алгоритмы и методы, заложенные в основу каждого модуля.

% Во второй главе описан разработанный пакет программ, состоящий из трёх частей: вычислительная, визуальная и аналитическая.
% Так как программ много, в первую очередь описываются потоки данных между ними, а также даётся краткое описание каждой из них.
% Далее подробно описывается каждая часть.
% Представлены вычислительные программы, их модули и методы, применяемые в них.
% После этого описываются программы для визуализации частиц и сетки, возникшие проблемы при их использовании, 
% повлиявшие на дальнейшее развитие проекта.
% В конце главы представлены программы для анализа результатов моделирования.
% Они сравнительно простые, но необходимы, так как используются для наиболее часто встречающихся задач анализа динамики жидкости.

В третьей главе приведены результаты тестирования разработанной модели SPH.
%
Представлена оценка эффективности распараллеливания вычислительных модулей и
сравнение максимального ускорения при использовании расчётов на CPU и на GPU.
%
Проведён анализ погрешности в дисперсионном соотношении численной модели.
%
Проведён анализ затухания волн на поверхности воды, рассчитана на её основне эффективная вязкость численной модели.
%
Проведено тестирование возможностей разработанного ПО для решения задач трёхмерного моделирования движения жидкости.
%
Результаты демонстрируют соответствие модели теоретическим ожиданиям и 
её применимость для широкого класса гидродинамических задач.


% Полный объём диссертации составляет
% \formbytotal{TotPages}{страниц}{у}{ы}{}, включая
% \formbytotal{totalcount@figure}{рисун}{ок}{ка}{ков} и
% \formbytotal{totalcount@table}{таблиц}{у}{ы}{}.
% Список литературы содержит
% \formbytotal{citenum}{наименован}{ие}{ия}{ий}.
